\section{Resulting Issues}
These problems can lead to several critical issues in educational applications of
general-purpose LLMs. The following consequences undermine the reliability and effectiveness
of such systems in supporting student learning:

\begin{itemize}
    \item Answers may include inaccuracies, as models draw upon broad but imperfect knowledge
          bases without verification against authoritative sources or course-specific materials.
    \item Some explanations may be irrelevant or missing key elements, failing to address
          the precise conceptual gaps or learning objectives relevant to a given learner.
    \item Answers may ignore details the instructor has already explained in class, resulting
          in misalignment between the model's output and the established curriculum narrative.
    \item Responses can sometimes contain hallucinations—plausible yet factually incorrect
          information—which poses a particular risk in high-stakes educational contexts where
          students may internalize erroneous concepts without critical scrutiny.
    \item The model fails to identify and explain prerequisite knowledge and fill background gaps.
          It does not determine the prerequisite knowledge gaps of students and
          attempts to explain general prerequisites, which are often inaccurate for any given
          individual.
\end{itemize}

These limitations collectively highlight the inadequacy of generic LLMs for nuanced pedagogical support. Without mechanisms for personalization, curricular alignment, and factual grounding, their deployment risks not only diminishing learning outcomes but also eroding trust in AI-assisted education. Addressing these challenges necessitates the development of more sophisticated, context-aware AI systems capable of adapting to individual learners, specific course content, and institutional requirements.
